\documentclass[11pt]{article}
\usepackage[margin=1in]{geometry}
\usepackage{amsmath}
\usepackage{booktabs}
\usepackage{enumitem}
\usepackage{hyperref}

\title{\textbf{ADR-02: Data Generation Strategy for Customer Query Analytics}}
\author{Architecture Decision Record}
\date{\today}

\begin{document}

\maketitle

\section{Status}
\textbf{Accepted}

\section{Context}

Customer query analytics requires high-quality annotated datasets to develop reliable classification systems and assess inter-annotator agreement. The project faces constraints around data availability, privacy regulations, and the need for controlled experimentation.

\subsection{Data Requirements}
\begin{itemize}
\item Realistic customer support ticket content across multiple difficulty levels
\item Multiple annotator perspectives to assess inter-annotator agreement
\item Hierarchical labels (L1: Technical/Billing/Account, L2: Specific subtypes)
\item Sufficient volume for statistical analysis and pattern detection
\item Diverse scenarios representing real-world annotation challenges
\end{itemize}

\subsection{Alternative Approaches Considered}
\begin{enumerate}
\item \textbf{Real Customer Data}: Use actual support tickets with anonymization
\item \textbf{Public Datasets}: Leverage existing customer service datasets
\item \textbf{Crowdsourced Annotation}: Hire external annotators for real data labeling
\item \textbf{Synthetic Data Generation}: AI-generated tickets with simulated annotator responses
\item \textbf{Hybrid Approach}: Combination of real data samples with synthetic augmentation
\end{enumerate}

\section{Decision}

We selected \textbf{LLM Based Synthetic Data Generation with Multi-Annotator Simulation} using the following architecture:

\subsection{Core Strategy}
\begin{itemize}
\item Generate 50 hyper-realistic customer support tickets across three difficulty tiers
\item Simulate 7 distinct annotator personas with varying skill levels and biases
\item Implement controlled policy ambiguities to create realistic disagreement patterns
\item Use tiered complexity (Easy: 15, Medium: 15, Hard: 20 samples) for comprehensive analysis
\end{itemize}

\subsection{Annotator Simulation Framework}
\begin{itemize}
\item \textbf{Tier 1 - Expert Annotators}: alex-001, beth-002, carlos-003 (high accuracy)
\item \textbf{Tier 2 - Medium Annotators}: diane-004, eric-005 (policy-confused but well-intentioned)  
\item \textbf{Tier 3 - Poor Annotators}: fatima-006, george-007 (low-effort, keyword-based)
\end{itemize}

\subsection{Complexity Stratification}
\begin{itemize}
\item \textbf{Easy Tickets (15)}: Clear, unambiguous requests with high expected agreement
\item \textbf{Medium Tickets (15)}: Policy ambiguities testing L1/L2 boundary decisions
\item \textbf{Hard Tickets (20)}: Stream-of-consciousness, emotionally charged, multi-issue scenarios
\end{itemize}

\section{Rationale}

\subsection{Advantages of Synthetic Approach}
\begin{itemize}
\item \textbf{Privacy Compliance}: Zero risk of customer data exposure or regulatory violations
\item \textbf{Controlled Experimentation}: Precisely engineered scenarios to test specific analytical capabilities
\item \textbf{Annotator Bias Simulation}: Realistic human annotation patterns without human resource costs
\item \textbf{Rapid Iteration}: Ability to generate additional samples or modify scenarios as needed
\item \textbf{Ground Truth Control}: Known complexity levels enable validation of analytical methods
\end{itemize}

\subsection{Technical Implementation Benefits}
\begin{itemize}
\item \textbf{Consistent Quality}: Every ticket engineered to serve specific analytical purposes
\item \textbf{Realistic Disagreement}: Annotator personas create authentic inter-annotator agreement patterns
\item \textbf{Scalable Generation}: Framework can produce additional datasets for future analysis
\item \textbf{Policy Testing}: Deliberate ambiguities validate annotation policy effectiveness
\end{itemize}

\section{Consequences}

\subsection{Positive Outcomes}
\begin{itemize}
\item \textbf{Risk Mitigation}: Eliminated data privacy and compliance concerns
\item \textbf{Analytical Validity}: Controlled dataset enables rigorous testing of IAA methods
\item \textbf{Cost Efficiency}: No annotation labor costs or real data acquisition expenses
\item \textbf{Reproducibility}: Consistent dataset enables reliable algorithm comparison and validation
\item \textbf{Educational Value}: Clear demonstration of annotation challenges and solution approaches
\end{itemize}

\subsection{Trade-offs and Limitations}
\begin{itemize}
\item \textbf{Synthetic Bias}: Generated content may miss nuances of real customer language
\item \textbf{Limited Vocabulary}: Constrained to AI-generated linguistic patterns
\item \textbf{Scale Constraints}: 50 samples sufficient for proof-of-concept but limited for production training
\item \textbf{Domain Gaps}: May not capture industry-specific terminology or edge cases
\end{itemize}

\subsection{Rejected Alternatives}
\begin{itemize}
\item \textbf{Real Customer Data}: Legal and privacy risks outweighed analytical benefits
\item \textbf{Public Datasets}: Available datasets lacked hierarchical structure and multi-annotator perspective
\item \textbf{Crowdsourced Annotation}: High cost and quality control challenges
\item \textbf{Hybrid Approach}: Complexity overhead without sufficient additional benefit
\end{itemize}

\section{Implementation Details}

\subsection{Data Generation Parameters}
\begin{itemize}
\item \textbf{Ticket Length}: 2-3 sentences (easy), 3-5 sentences (medium), 4-7 sentences (hard)
\item \textbf{Emotional Range}: Neutral/polite (easy) to highly emotional/frustrated (hard)
\item \textbf{Ambiguity Injection}: Systematic testing of Login vs. Password Reset boundary cases
\item \textbf{Realism Features}: Typos, grammar errors, stream-of-consciousness patterns in hard tier
\end{itemize}

\subsection{Quality Validation}
\begin{itemize}
\item Manual review of generated samples for realism and diversity
\item Verification that annotator personas produce expected disagreement patterns
\item Confirmation that difficulty tiers create appropriate analytical challenges
\item Cross-validation with subject matter experts on ticket authenticity
\end{itemize}

\subsection{Success Metrics}
\begin{itemize}
\item Inter-annotator agreement ranges: Easy ($\alpha > 0.8$), Medium ($0.4 < \alpha < 0.8$), Hard ($\alpha < 0.6$)
\item Annotator performance differentiation clearly visible in confusion matrices
\item Sufficient linguistic diversity for meaningful topic modeling and entity recognition
\item Realistic sentiment distribution across emotional spectrum
\end{itemize}

\section{Future Considerations}

\subsection{Evolution Path}
\begin{itemize}
\item \textbf{Scale Expansion}: Generate larger datasets (500+ samples) for production model training
\item \textbf{Real Data Integration}: Incorporate anonymized real samples for validation
\item \textbf{Domain Extension}: Adapt generation framework for other business contexts
\item \textbf{Dynamic Generation}: Real-time synthetic data creation for continuous testing
\end{itemize}

\subsection{Validation Requirements}
Quarterly assessment of synthetic data quality through:
\begin{itemize}
\item Comparison with emerging real customer ticket patterns
\item Validation against production model performance
\item Expert review of linguistic authenticity and business relevance
\item Extension testing to additional customer service scenarios
\end{itemize}

\end{document}